\documentclass[../../main.tex]{subfiles}
\begin{document}
\section{Route C, exact cases: one dimension, products, and the log-concave Dirichlet family}
\label{sec:cmh-exact-cases}

Section~\ref{sec:cmh-normalization} fixed the CMH estimate and proved that it dominates the
affine Poincar\'e constant. This section computes it exactly in the three classes where it is
tractable. Two of them --- the line and products --- are the expected calibrations, and they
already pin the constant: $\mathrm{CMH}(4)$ holds there and no smaller universal constant is
possible. The third is the first genuinely nonproduct family for which the route has a theorem
at all: every log-concave Dirichlet law satisfies $\mathrm{CMH}(4)$
(Theorem~\ref{thm:cmh-dirichlet}). Its proof is a homogeneous lift to independent Gamma
variables followed by a sharp Hessian-row minimization; in the Dirichlet argument the
log-concavity hypothesis is consumed in the final scalar angular minimization.

The section closes at the exact product endpoint: centered one-sided exponentials saturate
$\mathrm{CMH}(4)$ with \emph{zero} slack. Whether a log-concave perturbation raises the full CMH
Rayleigh quotient is a separate, open second-variation problem; changing its solenoidal term
alone is not decisive (Remark~\ref{rem:cmh-saturation-risk},
Question~\ref{q:cmh-solenoidal-perturbation}).

\subsection{The line}
\label{subsec:cmh-1d}

Let $\mu(\dd x)=\rho(x)\dd x$ be centered on an interval $(\ell,r)$ with variance $\sigma^2$.
Its canonical Stein kernel is the unique zero-flux solution of $(\tau\rho)'=-x\rho$,
\begin{equation}
  \tau(x)=-\frac1{\rho(x)}\int_\ell^xt\rho(t)\dd t,
  \label{eq:cmh-1d-stein}
\end{equation}
and it coincides with the moment-map Hessian in target coordinates. The Stein generator is
$L_\mu f=\tau f''-xf'=\rho^{-1}(\tau\rho f')'$, so that $h=-L_\mu f$ has
\begin{equation}
  \tau(x)f'(x)=-\frac1{\rho(x)}\int_\ell^xh(t)\rho(t)\dd t.
  \label{eq:cmh-1d-flux}
\end{equation}

\begin{theorem}[Exact one-dimensional CMH constant]
\label{thm:cmh-1d}
For every centered one-dimensional law for which the operators above are defined,
\begin{equation}
  \CMH(\mu)=\frac{\CP(\mu)}{\Var(\mu)}.
  \label{eq:cmh-1d-exact}
\end{equation}
In particular every one-dimensional log-concave law satisfies $\mathrm{CMH}(4)$, and the
constant $4$ cannot be lowered.
\end{theorem}

\begin{proof}
On the centered subspace let $D=\dd/\dd x$ and let $D_\mu^*u=-\rho^{-1}(\rho u)'$ be its
weighted adjoint, so that the inverse norm of the nonnegative Langevin operator $D_\mu^*D$ is
$\CP(\mu)$. If $v=(D_\mu^*D)^{-1}h$, then $Dv$ is the zero-boundary-flux solution of
$D_\mu^*u=h$; by \eqref{eq:cmh-1d-flux} it is the field $u=\tau f'$. Therefore
\[
  \norm u_2^2=\norm{Dv}_2^2
  =\inner{v}{D_\mu^*Dv}
  =\inner{h}{(D_\mu^*D)^{-1}h},
  \qquad
  \sup_{h\perp1}\frac{\norm u_2^2}{\norm h_2^2}=\CP(\mu).
\]
The CMH numerator is $\norm u_2^2/\sigma^2$ and its denominator is $\norm h_2^2$, which gives
\eqref{eq:cmh-1d-exact}. The one-dimensional specialization of the
Kannan--Lov\'asz--Simonovits bound \cite{KannanLovaszSimonovits1995}, recorded for example in
\cite[Eq.~(2.25)]{cattiaux2018poincare} with that attribution, is
$\CP(\mu)\le4\Var(\mu)$ and hence gives $\mathrm{CMH}(4)$.

For sharpness, let $X=Y-1$ with $Y\sim\mathrm{Exp}(1)$, so that $X$ is centered, has density
$e^{-(x+1)}\one_{x\ge-1}\dd x$, and has variance $1$. For $0<a<\tfrac12$, the test function
$f_a(x)=e^{a(x+1)}$ satisfies
\[
  \frac{\Var(f_a(X))}{\E\abs{f_a'(X)}^2}
  =\frac{(1-2a)^{-1}-(1-a)^{-2}}{a^2(1-2a)^{-1}}
  =\frac1{(1-a)^2}\longrightarrow4.
\]
Thus $\CP(X)=4$ by the preceding upper bound, and \eqref{eq:cmh-1d-exact} gives $\CMH(X)=4$.
\end{proof}

Theorem~\ref{thm:cmh-1d} is an \emph{identity}, not an inequality: on the line CMH carries
exactly the information of the affine Poincar\'e inequality and nothing more. This is
Remark~\ref{rem:cmh-stronger-than-kls} seen from the other side --- the solenoidal channel is
empty in dimension one.

\subsection{Products}
\label{subsec:cmh-product}

Let $\mu=\bigotimes_{i=1}^m\mu_i$ be a product of centered factors for which the canonical CMH
data are defined. Then $\Sigma=\diag(\Sigma_1,\dots,\Sigma_m)$,
$H=\diag(H_1,\dots,H_m)$, and $L_\mu=\sum_iL_i$, where the $L_i$ are commuting nonpositive
generators acting in their respective blocks.

\begin{theorem}[Product formula]
\label{thm:cmh-product}
For a product of centered CMH factors,
\begin{equation}
  \CMH\Bigl(\bigotimes_i\mu_i\Bigr)=\max_i\CMH(\mu_i).
  \label{eq:cmh-product}
\end{equation}
In particular, for centered one-dimensional factors this is
$\max_i\CP(\mu_i)/\Var(\mu_i)$. Consequently $\mathrm{CMH}(4)$ holds for every product of
one-dimensional log-concave laws, for block products of $\mathrm{CMH}(4)$ factors, and for their
invertible affine images.
\end{theorem}

\begin{proof}
Write $\nabla_i$ for the gradient in the $i$th block. Conditioning on all other blocks and
using the definition of $\CMH(\mu_i)$ gives
\[
  \E\inner{H_i\nabla_i g}{\Sigma_i^{-1}H_i\nabla_i g}
  \le \CMH(\mu_i)\,\E(L_i g)^2.
\]
For $i\ne j$, two integrations by parts give the block cross-term identity
\begin{equation}
  \E[(L_ig)(L_jg)]
  =\E\Tr\!\left(H_iD^2_{ij}g\,H_jD^2_{ji}g\right)
  =\E\norm{H_i^{1/2}D^2_{ij}g\,H_j^{1/2}}_{\HS}^2\ge0,
  \label{eq:cmh-product-cross}
\end{equation}
whence $\sum_i\E(L_ig)^2\le\E\bigl(\sum_iL_ig\bigr)^2=\E(L_\mu g)^2$. Summing the conditional
estimates proves ``$\le$'' in \eqref{eq:cmh-product}; testing on functions of a single block
proves ``$\ge$''. The one-dimensional formula now follows from
Theorem~\ref{thm:cmh-1d}, and the final affine-image claim uses only the invertible affine
invariance proved in \S\ref{subsec:cmh-conventions}.
\end{proof}

\begin{corollary}[Poincar\'e closure under linear images and convolution]
\label{cor:cmh-linear-images}
If independent factors each satisfy the affine Poincar\'e inequality with constant $C$, then so
does every linear image of their product --- in particular every sum of independent such
factors.
\end{corollary}

\begin{proof}
For $Y=T(X_1,\dots,X_m)$ apply the product affine Poincar\'e inequality to $f\circ T$. The
gradient energy is $\nabla f^\top T\Cov(X)T^\top\nabla f=\inner{\Cov(Y)\nabla f}{\nabla f}$, and
$\Var_Yf=\Var_X(f\circ T)$.
\end{proof}

Note that Corollary~\ref{cor:cmh-linear-images} is a statement about $\CPaff$, which is stable
under noninvertible maps; the canonical CMH constant itself is claimed invariant only under
invertible affine maps (\S\ref{subsec:cmh-conventions}).

\subsection{The Dirichlet family: geometry and normalization}
\label{subsec:cmh-dirichlet-setup}

Let $P\sim\Dir(\alpha_1,\dots,\alpha_m)$ with $\alpha_i\ge1$, which is exactly the log-concave
parameter range, and put $A=\sum_i\alpha_i$, $q_i=\alpha_i/A$. Write
\begin{equation}
  C(p)=\diag(p)-pp^\top,
  \qquad
  L_\alpha g=\Tr\bigl(C(p)D^2g\bigr)+\inner{\alpha-Ap}{\nabla g}
  \label{eq:cmh-wright-fisher}
\end{equation}
for the Wright--Fisher generator; both expressions are independent of the ambient extension of
$g$ because $C(p)\one=0$ and $\sum_i(\alpha_i-Ap_i)=0$.

On $\R^m/\mathrm{span}\{\one\}$ consider
$\varphi(y)=A\log\bigl(\sum_ie^{y_i/A}\bigr)-q\cdot y$. Its gradient is $p-q$ with
$p_i=e^{y_i/A}/\sum_je^{y_j/A}$, its Hessian is $C(p)/A$, and the softmax Jacobian identifies
the pushforward of $e^{-\varphi}$ with the centered law $P-q$. Hence the canonical moment-map
data are
\begin{equation}
  H(p)=\tfrac1AC(p),\qquad L_\mu=\tfrac1AL_\alpha,
  \qquad
  \Sigma=\frac1{A(A+1)}\Bigl(\diag(\alpha)-\frac{\alpha\alpha^\top}A\Bigr).
  \label{eq:cmh-dirichlet-data}
\end{equation}
For every tangent vector $v$, meaning $\sum_iv_i=0$,
\begin{equation}
  v^\top\Sigma^\dagger v=A(A+1)\sum_i\frac{v_i^2}{\alpha_i},
  \label{eq:cmh-dirichlet-pseudoinverse}
\end{equation}
because $x=A(A+1)\diag(\alpha)^{-1}v$ satisfies $\Sigma x=v$ exactly, and the remaining
$\ker\Sigma=\mathrm{span}\{\one\}$ ambiguity pairs to zero against $v$.

Setting
\begin{equation}
  u_i(p)=\bigl(C(p)\nabla g(p)\bigr)_i,
  \qquad
  d_\alpha(g)=\E\sum_i\frac{u_i^2}{\alpha_i},
  \qquad
  n_\alpha(g)=\E(L_\alpha g)^2,
  \label{eq:cmh-dirichlet-dn}
\end{equation}
equations \eqref{eq:cmh-dirichlet-data}--\eqref{eq:cmh-dirichlet-pseudoinverse} turn
Definition~\ref{def:cmh} into the statement $A(A+1)d_\alpha(g)\le4n_\alpha(g)$.

\begin{theorem}[Log-concave Dirichlet CMH]
\label{thm:cmh-dirichlet}
For every $m\ge2$, every $\alpha_1,\dots,\alpha_m\ge1$, and every $g$ in the core,
\begin{equation}
  A(A+1)\,\E\sum_{i=1}^m
  \frac{\bigl(C(P)\nabla g(P)\bigr)_i^2}{\alpha_i}
  \le4\,\E\bigl(L_\alpha g(P)\bigr)^2.
  \label{eq:cmh-dirichlet-main}
\end{equation}
Equivalently, every log-concave Dirichlet law satisfies $\mathrm{CMH}(4)$.
\end{theorem}

\subsection{The independent Gamma lift}
\label{subsec:cmh-gamma-lift}

Let $Y_i\sim\GammaLaw(\alpha_i,1)$ be independent, $S=\sum_iY_i\sim\GammaLaw(A,1)$ and
$P_i=Y_i/S$; then $P\sim\Dir(\alpha)$ and $S\perp P$. Lift $g$ homogeneously of degree zero by
$G(Y)=g(Y/S)$. The product-Gamma (Laguerre) generator is
$\calL_\Gamma G=\sum_i\bigl(Y_iG_{ii}+(\alpha_i-Y_i)G_i\bigr)$, and its integrated Bochner
identity reads
\begin{equation}
  N_\Gamma(G):=\E(\calL_\Gamma G)^2
  =\E\Bigl[\sum_iY_iG_i^2+\sum_{i,j}Y_iY_jG_{ij}^2\Bigr].
  \label{eq:cmh-gamma-bochner}
\end{equation}
Put $D_i(G)=\E[Y_i^2G_i^2]/\alpha_i$ and $D_\Gamma(G)=\sum_iD_i(G)$.

\begin{lemma}[Exact Gamma row completion]
\label{lem:cmh-gamma-completion}
For $\alpha_i\ge1$,
\begin{equation}
  N_\Gamma(G)-\tfrac14D_\Gamma(G)=\sum_i\E R_i+\sum_i\delta_iD_i(G),
  \label{eq:cmh-gamma-completion}
\end{equation}
where
\begin{equation}
  R_i=Y_i^2\Bigl\lvert G_{ii}-\frac{G_i}{\alpha_i+1}\Bigr\rvert^2
      +\sum_{j\ne i}Y_iY_j\abs{G_{ij}}^2\ \ge0,
  \qquad
  \delta_i=\frac{\alpha_i^2}{(\alpha_i+1)^2}-\frac14\ \ge0.
  \label{eq:cmh-Ri-deltai}
\end{equation}
\end{lemma}

\begin{proof}
For $Y\sim\GammaLaw(a,1)$, integrating $\E[Y^2(u^2)']$ by parts against $y^{a-1}e^{-y}$ gives
$\E[Y^2uu']=-\tfrac12\E[((a+1)Y-Y^2)u^2]$. Apply this with $u=G_i$, $u'=G_{ii}$, $a=\alpha_i$
and expand the first square in \eqref{eq:cmh-Ri-deltai}: the cross term equals
$\E[Y_iG_i^2]-\E[Y_i^2G_i^2]/(\alpha_i+1)$, so
\[
  \E R_i=\E\Bigl[Y_iG_i^2+\sum_jY_iY_jG_{ij}^2\Bigr]
        -\frac{\alpha_i}{(\alpha_i+1)^2}\E[Y_i^2G_i^2],
\]
using $\tfrac1{(\alpha_i+1)^2}-\tfrac1{\alpha_i+1}=-\tfrac{\alpha_i}{(\alpha_i+1)^2}$. Summing
over $i$ recovers $N_\Gamma(G)$ from \eqref{eq:cmh-gamma-bochner} and leaves
$-\sum_i\alpha_i^2(\alpha_i+1)^{-2}D_i(G)$; subtracting $D_\Gamma/4$ gives
\eqref{eq:cmh-gamma-completion}. Finally $\delta_i\ge0$ is equivalent to
$\alpha_i/(\alpha_i+1)\ge\tfrac12$, i.e.\ to $\alpha_i\ge1$.
\end{proof}

\begin{remark}[Where log-concavity is spent]
\label{rem:cmh-dirichlet-logconcavity}
Lemma~\ref{lem:cmh-gamma-completion} already proves $\mathrm{CMH}(4)$ for the product-Gamma law,
since $R_i\ge0$ and $\delta_i\ge0$; for that product-Gamma consequence, $\alpha_i\ge1$ is used
only in the sign of $\delta_i$. The Dirichlet proof does not discard $\delta_i$: it retains its
exact value inside $F_A(\alpha_i,P_i)$. There the log-concavity hypothesis is consumed when
Lemma~\ref{lem:cmh-angular-coefficient} is applied with $a=\alpha_i\ge1$.
\end{remark}

\begin{lemma}[Hessian-row minimization]
\label{lem:cmh-row-min}
With $G$ homogeneous of degree zero,
\begin{equation}
  R_i\ \ge\ \frac{Y_i}S\Bigl(1+\frac{Y_i}{\alpha_i+1}\Bigr)^2\abs{G_i}^2 .
  \label{eq:cmh-row-min}
\end{equation}
\end{lemma}

\begin{proof}
Euler's identity $\sum_jY_jG_j=0$, differentiated in $Y_i$, gives the single linear constraint
$\sum_jY_jG_{ij}=-G_i$. Set $t_i=Y_i\bigl(G_{ii}-G_i/(\alpha_i+1)\bigr)$ and $t_j=Y_jG_{ij}$
for $j\ne i$, so that
$\sum_jt_j=-G_i\bigl(1+Y_i/(\alpha_i+1)\bigr)$ and
$R_i=t_i^2+\sum_{j\ne i}(Y_i/Y_j)t_j^2$. The reciprocals of these quadratic weights sum to
$1+\sum_{j\ne i}Y_j/Y_i=S/Y_i$, and the minimum of $\sum_jw_jt_j^2$ under $\sum_jt_j=c$ is
$c^2/\sum_jw_j^{-1}$, which is the right-hand side of \eqref{eq:cmh-row-min}.
\end{proof}

Homogeneity also gives the exact dictionary between the lift and the simplex:
$\calL_\Gamma G=S^{-1}L_\alpha g$ and $Y_iG_i=u_i(P)$. With $A>2$ one has $\E S^{-1}=(A-1)^{-1}$
and $\E S^{-2}=z_A^{-1}$ where
\begin{equation}
  z_A=(A-1)(A-2),
  \label{eq:cmh-zA}
\end{equation}
so independence of $S$ and $P$ yields
\begin{equation}
  N_\Gamma(G)=\frac{n_\alpha(g)}{z_A},
  \qquad
  D_\Gamma(G)=d_\alpha(g),
  \label{eq:cmh-radial-relations}
\end{equation}
and, substituting $Y_i=SP_i$ and $G_i=u_i/(SP_i)$ into \eqref{eq:cmh-row-min},
\begin{equation}
  \E R_i\ \ge\ \E_P\,u_i^2
  \Bigl[\frac1{z_AP_i}+\frac2{(A-1)(\alpha_i+1)}+\frac{P_i}{(\alpha_i+1)^2}\Bigr].
  \label{eq:cmh-integrated-row}
\end{equation}

\subsection{The scalar angular inequality and the proof}
\label{subsec:cmh-dirichlet-proof}

\begin{lemma}[Angular coefficient]
\label{lem:cmh-angular-coefficient}
Let $A\ge3$, $a\ge1$, $0<p\le1$, $z=(A-1)(A-2)$, and
\begin{equation}
  F_A(a,p)=\frac ap+\frac{2a(A-2)}{a+1}+\frac{za(a+p)}{(a+1)^2}-\frac z4 .
  \label{eq:cmh-FA}
\end{equation}
Then
\begin{equation}
  F_A(a,p)\ \ge\ A-\tfrac12+s_A,
  \qquad
  s_A=\begin{cases}\dfrac{z-2}4,&z\le4,\\[6pt]\sqrt z-\dfrac32,&z\ge4,\end{cases}
  \label{eq:cmh-FA-min}
\end{equation}
and $s_A\ge0$ for $A\ge3$.
\end{lemma}

\begin{proof}
For fixed $a$ the $p$-dependent part is $a\bigl(p^{-1}+zp(a+1)^{-2}\bigr)$, minimized on
$(0,1]$ at $p=1$ when $\sqrt z\le a+1$ and at $p_*=(a+1)/\sqrt z$ otherwise.

If $z\le4$ then $\sqrt z\le2\le a+1$ for all $a\ge1$, so $p=1$ and
$F_A=a+2a(A-2)/(a+1)+za/(a+1)-z/4$, each summand increasing in $a$; the minimum at $a=1$ is
$A-1+z/4=A-\tfrac12+(z-2)/4$.

If $z\ge4$, then at $p_*$ the two $p$-terms are equal and sum to $2a\sqrt z/(a+1)$, so with
$r=a/(a+1)$ the value is $2r(\sqrt z+A-2)+zr^2-z/4$, increasing in $r$; since $r\ge\tfrac12$ for
$a\ge1$ the minimum is at $a=1$, $r=\tfrac12$, and equals $A-2+\sqrt z=A-\tfrac12+\sqrt z-\tfrac32$.
In the complementary region the expression is increasing in $a$ with a boundary value no
smaller. Finally $A\ge3$ gives $z\ge2$, so $s_A\ge0$ in both branches.
\end{proof}

\begin{proof}[Proof of Theorem~\ref{thm:cmh-dirichlet}]
If $m=2$ and $A<3$, the law is one dimensional and Theorem~\ref{thm:cmh-1d} applies. In every
remaining case $A\ge3$ (automatically when $m\ge3$), so $z_A\ge2>0$ and the inverse moments
$\E S^{-1},\E S^{-2}$ used in \eqref{eq:cmh-radial-relations} are finite.

Combine \eqref{eq:cmh-gamma-completion}, \eqref{eq:cmh-Ri-deltai} and
\eqref{eq:cmh-integrated-row}. Since $D_i(G)=\E[u_i^2]/\alpha_i$, the total coefficient of
$\E_P[u_i^2]$ in $N_\Gamma(G)-\tfrac14D_\Gamma(G)$ is
\[
  \frac1{z_AP_i}+\frac2{(A-1)(\alpha_i+1)}+\frac{P_i}{(\alpha_i+1)^2}+\frac{\delta_i}{\alpha_i},
\]
and multiplying it by $z_A\alpha_i$ gives exactly $F_A(\alpha_i,P_i)$ of \eqref{eq:cmh-FA},
using $z_A/(A-1)=A-2$. By Lemma~\ref{lem:cmh-angular-coefficient},
\[
  N_\Gamma(G)-\tfrac14D_\Gamma(G)
  \ \ge\ \frac{A-\tfrac12+s_A}{z_A}\sum_i\frac{\E_P[u_i^2]}{\alpha_i}
  =\frac{A-\tfrac12+s_A}{z_A}\,d_\alpha(g).
\]
Substituting \eqref{eq:cmh-radial-relations}, multiplying by $z_A$ and using the identity
$z_A/4+A-\tfrac12=A(A+1)/4$ gives
\begin{equation}
  n_\alpha(g)\ \ge\ \Bigl(\frac{A(A+1)}4+s_A\Bigr)d_\alpha(g),
  \label{eq:cmh-dirichlet-surplus}
\end{equation}
and $s_A\ge0$ yields \eqref{eq:cmh-dirichlet-main}.
\end{proof}

\begin{corollary}[Quantitative surplus]
\label{cor:cmh-dirichlet-surplus}
For $A>3$, $n_\alpha(g)\ge\bigl(\tfrac14+\eps_A\bigr)A(A+1)d_\alpha(g)$ with
$\eps_A=s_A/(A(A+1))>0$. Equivalently
$\CMH\le4\bigl(1+4s_A/(A(A+1))\bigr)^{-1}<4$: the constant $4$ is attained only in the limit.
\end{corollary}

\begin{proof}
When $A>3$ the Gamma branch of the preceding proof applies, including when $m=2$, and
\eqref{eq:cmh-dirichlet-surplus} gives the asserted estimate. Moreover $z_A>2$, so either
$s_A=(z_A-2)/4>0$ or $s_A=\sqrt{z_A}-3/2>0$.
\end{proof}

\begin{corollary}[Affine Poincar\'e consequences]
\label{cor:cmh-dirichlet-poincare}
Every log-concave Dirichlet law satisfies $\CPaff\le4$; so does every linear image, product, and
convolution of independent such laws.
\end{corollary}

\begin{proof}
Combine Theorems~\ref{thm:cmh-implies-affine-poincare} and \ref{thm:cmh-dirichlet} with
Corollary~\ref{cor:cmh-linear-images}.
\end{proof}

Qualitative dimension-free KLS bounds for simplices and conservative Gamma models are prior
art; see \cite[\S1.2]{KolesnikovMilman2016OrliczKLS} as a secondary pointer to that literature.
The specific content of this proof is the constant $4$, the quantitative surplus, and the
CMH-level estimate; no broader priority claim is intended.

\begin{remark}[Closure and aggregation]
\label{rem:cmh-dirichlet-closure}
The proof is written for smooth $g$ with controlled boundary behavior. The closure argument
takes polynomials on the simplex, lifts them after a radial cutoff $S\ge\eps$, applies the Gamma
identities and lets $\eps\downarrow0$; the required inverse Gamma moments are finite because
$A\ge3$ throughout the Gamma branch. Polynomials form a core for $L_\alpha$, so the closed forms
extend \eqref{eq:cmh-dirichlet-main} to $\Dom(L_\alpha)$; only the residual branch $m=2$, $A<3$
uses the one-dimensional no-flux closure instead.

The estimate is also compatible with Dirichlet aggregation. If
$Q\sim\Dir(\beta_1,\dots,\beta_N)$, the coordinates are partitioned into blocks $G_i$, and
$P_i=\sum_{a\in G_i}Q_a$, $\alpha_i=\sum_{a\in G_i}\beta_a$, then for functions of $P$ alone the
generator is the coarse one and
$\E\bigl[\sum_{a\in G_i}Q_a^2/\beta_a\mid P\bigr]
 =P_i^2(\alpha_i+\abs{G_i})/(\alpha_i(\alpha_i+1))\ge P_i^2/\alpha_i$,
so CMH for the fine law implies CMH for its aggregation. This is a consistency check on
Theorem~\ref{thm:cmh-dirichlet}, not an ingredient of it.
\end{remark}

\begin{remark}[Sharpness on simplex faces]
\label{rem:cmh-dirichlet-sharp}
The constant $4$ cannot be improved uniformly over the family. The face law
$\mathrm{Beta}(1,b)$ is one dimensional, so Theorem~\ref{thm:cmh-1d} applies; its Neumann
spectral problem reduces to Bessel's equation and gives
$\CMH(\mathrm{Beta}(1,b))=(b+1)^2(b+2)/\bigl(b\,j_{b/2,1}^2\bigr)$ with $j_{\nu,1}$ the first
positive zero of $J_\nu$. Since $j_{b/2,1}\sim b/2$, this tends to $4$; after rescaling
$\mathrm{Beta}(1,b)$ converges to the one-sided exponential. The extremal mechanism on the
simplex is therefore the same one-dimensional exponential boundary mechanism that fixes the
constant in Theorem~\ref{thm:cmh-1d}. Corollary~\ref{cor:cmh-dirichlet-surplus} says the
interior of the family is strictly safer.
\end{remark}

\subsection{The saturation risk and the decisive test}
\label{subsec:cmh-saturation}

\begin{remark}[Exact product saturation and an open perturbative test]
\label{rem:cmh-saturation-risk}
By Theorems~\ref{thm:cmh-1d} and \ref{thm:cmh-product}, a product of centered one-sided
exponentials has $\CMH=4$ \emph{exactly}. Proposition~\ref{prop:cmh-hodge} gives, for every
admissible test function, the exact splitting of the CMH numerator into its affine Poincar\'e
part and the nonnegative solenoidal term $\E\inner{w}{\Sigma^{-1}w}$.

No perturbative conclusion follows from those two facts alone. Under a perturbation the
covariance and its inverse, the canonical Stein kernel, both numerator channels, the
denominator, and the optimizing test function may all vary. In particular, an increase of the
solenoidal term for one test function does not by itself imply an increase of the full CMH
Rayleigh quotient. Whether an admissible perturbation raises $\CMH$ above $4$ is the explicit
open second-variation problem in Question~\ref{q:cmh-solenoidal-perturbation}; a positive answer
would refute $\mathrm{CMH}(4)$ but would not by itself refute Conjecture~\ref{conj:kls}.
\end{remark}

\begin{question}[Second-order solenoidal excess at the product endpoint]
\label{q:cmh-solenoidal-perturbation}
Take the product moment potential $\psi_0(s,t)=\phi(s)+t^2/2$ with $\phi$ the one-sided
exponential moment potential, and perturb it by $\psi_\eps=\psi_0+\eps\,a(s)b(t)$ within the
smooth strictly convex class. Compute the second variation of $\CMH$ at $\eps=0$, splitting the
CMH numerator through Proposition~\ref{prop:cmh-hodge} into its gradient and solenoidal parts.
Either exhibit an admissible perturbation with strictly positive second variation --- refuting
$\mathrm{CMH}(4)$ and forcing the route to a larger constant --- or identify the structural
reason the solenoidal part vanishes to second order at a saturating product.
\end{question}

Question~\ref{q:cmh-solenoidal-perturbation} is the sharpest available probe of the route
because it attacks the endpoint rather than the machinery, and because both ingredients are
already exact: the saturation value comes from Theorem~\ref{thm:cmh-product} and the splitting
from Proposition~\ref{prop:cmh-hodge}. Together with Question~\ref{q:gate-zero} it forms the
falsification layer of Route C; Questions~\ref{q:mm-invariant-lift} and
\ref{q:mm-square-root-commutator} remain the construction layer.
\end{document}
